\section{The base language}

To show the relation among the clairvoyance semantics, the demand semantics, and
functional logic programming, we first define a base language $B$ that can be
translated into those semantics.

We show the syntax of $B$ in \cref{fig:syntax-base}. The base language is a
simple calculus with booleans, lists, and lazy thunks. Its term contains
variables, let-bindings, conditionals, recursions based on $\tfoldrZ$. In
addition, it contains the explicit \( \tlazyZ \) and \( \tforceZ \) for wrapping
a computation in a lazy thunk and forcing a computation in a thunk. This
decouples our analysis from any particular evaluation model. Finally, we have an
explicit \( \ttickZ \) term to indicate incurring one unit of computation cost.

The language does not have functions or general recursions. This is because the
demand semantics does not support those
features~\cite{bjerner1989,demand-semantics}.\ly{But functional logic
  programming does not support this in general, right? (Perhaps with the
  exception of Verse.)}

\begin{figure}[t]
\begin{center}
\begin{bnf}
  \( A \) ::= \TBool{} // \TList{A} // \TThunk{A} ;;

  $t$ ::= $x$ // \tlet{x}{t}{t} | \ttrue{} // \tfalse{} // \tif{t}{t}{t} |
  \tnil{} // \tcons{t}{t} // \tfoldr{x}{x}{t}{t}{t} | \tlazy{t} // \tforce{t} //
  \ttick{t}
\end{bnf}
\end{center}
\caption{The syntax of our base language B.}\label{fig:syntax-base}
\end{figure}

\begin{figure}[t]
  \begin{center}
    \begin{mathpar}
      \inferrule{\Gamma \vdash^\LB t_1 : A \\ \Gamma \vdash^\LB t_2 :
        \TThunk{(\TList{A})}}{\Gamma \vdash^\LB \tcons{t_1}{t_2} : \TList{A}} \\
      \inferrule{\Gamma, x_1 : A, x_2 : \TThunk{B} \vdash^\LB t_1 : B \\ \Gamma
        \vdash^\LB t_2 : B \\ \Gamma \vdash^\LB t_3 : \TList{A}}{\Gamma \vdash^\LB \tfoldr{x_1}{x_2}{t_1}{t_2}{t_3} : B} \\
      \inferrule{\Gamma \vdash^\LB t : A}{\Gamma \vdash^\LB \tlazy{t} : \TThunk{A}} \qquad
      \inferrule{\Gamma \vdash^\LB t : \TThunk{A}}{\Gamma \vdash^\LB \tforce{t} : A} \qquad
      \inferrule{\Gamma \vdash^\LB t : A}{\Gamma \vdash^\LB \ttick{t} : A}
    \end{mathpar}
  \end{center}
  \caption{Typing rules in $B$ that are related to lazy thunks and cost.}\label{fig:typing-base}
\end{figure}

Most of the typing rules in $B$ are standard, so we only show the typing rules
that are related to lazy thunks and the computation cost in
\cref{fig:typing-base}. Notably, the cons of lists $t_1 :: t_2$ always takes a
$t_1$ of type $A$ and a $t_2$ of type $\TThunk{(\TList{A})}$, forcing the tail
of a list to be a lazy thunk.\ly{But why?} Similarly, the function in $\tfoldrZ$
always assumes that the second paramter is a thunk.\ly{Why?}

\paragraph{Exact values and approximations.}

Each type \( A \) comes equipped with a set \( \Eval{A} \) of \emph{exact
  values} and a set \( \Approx{A} \) of \emph{approximation values}, as defined
in \cref{fig:exact-approx}. The set of exact values \( \Eval{A} \) represents
all possible values of $A$ that contain \emph{no} thunk. The set of
approximation values \( \Approx{A} \), on the other hand, represents\ly{How to
  explain the set of approximation values?}

\begin{figure}[t]
\begin{center}
  \begin{align*}
    \Eval{\TBool} &= \SBool \\
    \Eval{\TThunk{A}} &= \Eval{A} \\
    \Eval{\TList{A}} &= \braces{\snil} \cup \setbuilder{\scons{v_1}{v_2}}{v_1 \in \Eval{A}, v_2 \in \Eval{\TList{A}}} \\ \\
    \Approx{\TBool} &= \SBool \\
    \Approx{\TThunk{A}} &= \set{\sundefined} \cup \setbuilder{\sthunk{v}}{v \in \Approx{A}} \\
    \Approx{\TList{A}} &= \set{\snil} \cup \setbuilder{\scons{v_1}{v_2}}{v_1 \in \Approx{A}, v_2 \in \Approx{\TThunk{(\TList{A})}}}
  \end{align*}
\end{center}
\caption{Exact values and approximation values.\ly{What is the type of
    \( \Eval{A} \) and \( \Approx{A} \)?}}\label{fig:exact-approx}
\end{figure}

\subsection{Evaluation semantics}

Evaluation semantics\ly{Why does the evaluation semantics use the same notation
  as the set of exact values?}

\begin{align*}
  \EE{x}{\gamma} &= \gamma(x) \\
  \EE{\tlet{x}{t_1}{t_2}}{\gamma} &= \EE{t_2}{\gamma, x \mapsto \EE{t_1}{\gamma}} \\
  \EE{\ttrue}{\gamma} &= \strue \\
  \EE{\tfalse}{\gamma} &= \sfalse \\
  \EE{\tif{t_1}{t_2}{t_3}}{\gamma} &=
    \begin{cases}
      \EE{t_2}{\gamma} & \text{if \( \EE{t_{1}}{\gamma} = \strue \)} \\
      \EE{t_3}{\gamma} & \text{otherwise}
    \end{cases} \\
  \EE{\tnil}{\gamma} &= \varepsilon \\
  \EE{\tcons{t_1}{t_2}}{\gamma} &= \scons{\EE{t_1}{\gamma}}{\EE{t_2}{\gamma}} \\
  \EE{\tfoldr{x_1}{x_2}{t_1}{t_2}{t_3}}{\gamma} &= \sfoldr{(v_1, v_2) \mapsto \EE{t_1}{\gamma, x_1 \mapsto v_1, x_2 \mapsto v_2}}{\EE{t_2}{\gamma}}{\EE{t_3}{\gamma}} \\
  \EE{\tlazy{t}}{\gamma} &= \EE{t}{\gamma} \\
  \EE{\tforce{t}}{\gamma} &= \EE{t}{\gamma} \\
  \EE{\ttick{t}}{\gamma} &= \EE{t}{\gamma}
\end{align*}

\begin{align*}
  \sfoldr{f}{y}{(x_1, x_2, \ldots, x_n)} = f(x_1, f(x_2, \cdots f(x_n, y) \cdots))
\end{align*}

\begin{lemma}
  If \( \TB{\Gamma}{t}{A} \) and \( \gamma \in \Eval{\Gamma} \), then \( \EE{t}{\gamma} \in \Eval{A} \).
\end{lemma}

\subsection{Clairvoyance semantics}

Clairvoyance semantics has the judgment schema \( \EC{\gamma}{t}{c}{v} \).  Intuitively, \( c \) is the number of \( \texttt{tick} \)s encountered while evaluating \( t \) to \( v \).

\begin{center}
  \begin{prooftree}
    \hypo{\gamma(x) = v}
    \infer1{\EC{\gamma}{x}{0}{v}}
  \end{prooftree}
\end{center}

\begin{center}
  \begin{prooftree}
    \hypo{\EC{\gamma}{t_1}{c_1}{v_1}}
    \hypo{\EC{\gamma, x \mapsto v_1}{t_2}{c_2}{v_2}}
    \infer2{\EC{\gamma}{\tlet{x}{t_1}{t_2}}{c_1 + c_2}{v_2}}
  \end{prooftree}
\end{center}

\begin{center}
  \begin{prooftree}
    \infer0{\EC{\gamma}{\ttrue}{0}{\strue}}
  \end{prooftree}
\end{center}

\begin{center}
  \begin{prooftree}
    \infer0{\EC{\gamma}{\tfalse}{0}{\sfalse}}
  \end{prooftree}
\end{center}

\begin{center}
  \begin{prooftree}
    \hypo{\EC{\gamma}{t_1}{c_1}{\strue}}
    \hypo{\EC{\gamma}{t_2}{c_2}{v}}
    \infer2{\EC{\gamma}{\tif{t_1}{t_2}{t_3}}{c_1 + c_2}{v}}
  \end{prooftree}
\end{center}

\begin{center}
  \begin{prooftree}
    \hypo{\EC{\gamma}{t_1}{c_1}{\sfalse}}
    \hypo{\EC{\gamma}{t_3}{c_3}{v}}
    \infer2{\EC{\gamma}{\tif{t_1}{t_2}{t_3}}{c_1 + c_3}{v}}
  \end{prooftree}
\end{center}

\begin{center}
  \begin{prooftree}
    \infer0{\EC{\gamma}{\tnil}{0}{\snil}}
  \end{prooftree}
\end{center}

\begin{center}
  \begin{prooftree}
    \hypo{\EC{\gamma}{t_1}{c_1}{v_1}}
    \hypo{\EC{\gamma}{t_2}{c_1}{v_2}}
    \infer2{\EC{\gamma}{\tcons{t_1}{t_2}}{c_1 + c_2}{\scons{v_1}{v_2}}}
  \end{prooftree}
\end{center}

\begin{center}
  \begin{prooftree}
    \hypo{\EC{\gamma}{t_3}{c_1}{v_1}}
    \hypo{\ECF{\gamma}{x_1}{x_2}{t_1}{t_2}{v_1}{c_2}{v_2}}
    \infer2[\rulename{L-Foldr}]{\EC{\gamma}{\tfoldr{x_1}{x_2}{t_1}{t_2}{t_3}}{c_1 + c_2}{v_2}}
  \end{prooftree}
\end{center}

\begin{center}
  \begin{prooftree}
    \hypo{\EC{\gamma}{t}{c}{v}}
    \infer1[\rulename{L-Lazy-Thunk}]{\EC{\gamma}{\tlazy{t}}{c}{\sthunk{v}}}
  \end{prooftree}
\end{center}

\begin{center}
  \begin{prooftree}
    \infer0[\rulename{L-Lazy-Undefined}]{\EC{\gamma}{\tlazy{t}}{0}{\sundefined}}
  \end{prooftree}
\end{center}

\begin{center}
  \begin{prooftree}
    \hypo{\EC{\gamma}{t}{c}{\sthunk{v}}}
    \infer1[\rulename{L-Force}]{\EC{\gamma}{\tforce{t}}{c}{v}}
  \end{prooftree}
\end{center}

\begin{center}
  \begin{prooftree}
    \hypo{\EC{\gamma}{t}{c}{v}}
    \infer1{\EC{\gamma}{\ttick{t}}{c + 1}{v}}
  \end{prooftree}
\end{center}

The crucial rules are \rulename{L-Lazy-Thunk}, \rulename{L-Lazy-Undefined}, \rulename{L-Force}.  \rulename{L-Lazy-Thunk} and \rulename{L-Lazy-Undefined} say that, to evaluate a \( \tlazyZ \) term, we nondeterministically either evaluate its subterm and wrap it in a \( \sthunkZ \) or substitute an ``empty'' value \( \sundefined \).  \rulename{L-Force} says that we can only evaluate a value of \( \TThunkZ \) type if it is in fact a \( \sthunkZ \); in other words, if we need the result of a \( \tlazyZ \) term, then we must have actually evaluated it via \rulename{L-Lazy-Thunk}.

In the original formulation of clairvoyance semantics, nondeterminism was introduced at \( \tletZ \) bindings: when evaluating \( \tlet{x}{t_1}{t_2} \), either \( x \) was immediately bound to the result of evaluting \( t_1 \), or else \( x \) was not bound at all.  But our base language explicitly marks laziness points with \( \tthunkZ \) terms, so we must introduce nondeterminism there.

Note that \( \tfoldrZ \) is evaluated by first evaluting the list argument to a value and then iterating over that value.  The iteration is expressed in the following auxiliary relation.
\begin{center}
  \begin{prooftree}
    \hypo{\EC{\gamma}{t_2}{c}{v}}
    \infer1[\rulename{C-Foldr-Nil}]{\ECF{\gamma}{x_1}{x_2}{t_1}{t_2}{\snil}{c}{v}}
  \end{prooftree}
\end{center}
\begin{center}
  \begin{prooftree}
    \hypo{\EC{\gamma, x_1 \mapsto v_1, x_2 \mapsto \sundefined}{t_1}{c}{v}}
    \infer1[\rulename{C-Foldr-Undefined}]{\ECF{\gamma}{x_1}{x_2}{t_1}{t_2}{\scons{v_1}{\sundefined}}{c}{v}}
  \end{prooftree}
\end{center}
\begin{center}
  \begin{prooftree}
    \hypo{\ECF{\gamma}{x_1}{x_2}{t_1}{t_2}{v_2}{c_1}{v_3}}
    \hypo{\EC{\gamma, x_1 \mapsto v_1, x_2 \mapsto \sthunk{v_3}}{t_1}{c_2}{v}}
    \infer2[\rulename{C-Foldr-Thunk}]{\ECF{\gamma}{x_1}{x_2}{t_1}{t_2}{\scons{v_1}{\sthunk{v_2}}}{c_1 + c_2}{v}}
  \end{prooftree}
\end{center}

\begin{lemma}
  If \( \TB{\Gamma}{t}{A} \), \( \gamma \in \Approx{\Gamma} \), and \( \EC{\gamma}{t}{\relax}{v} \), then \( v \in \Approx{A} \).
\end{lemma}

\begin{lemma}
  If \( \TB{\Gamma}{t}{A} \) and \( \gamma \in \Approx{\Gamma} \), then \( \EC{\gamma}{t}{\relax}{\relax} \).
\end{lemma}

\subsection{Demand semantics}

\begin{align*}
  \ED{x}{\gamma}{a} &= (\braces{x \mapsto a}, 0) \\
  \ED{\ttrue}{\gamma}{a} &= (\bot, 0) \\
  \ED{\tfalse}{\gamma}{a} &= (\bot, 0) \\
  \ED{\tnil}{\gamma}{a} &= (\bot, 0) \\
  \ED{\tcons{t_1}{t_2}}{\gamma}{a} &= \ED{t_1}{\gamma}{a} \sqcup \ED{t_2}{\gamma}{a}
\end{align*}
